\documentclass[11pt]{article}
\usepackage[T1]{fontenc}
\usepackage[utf8]{inputenc}
\usepackage{lmodern}
\usepackage[margin=1in]{geometry}
\usepackage{microtype}
\usepackage{amsmath,amssymb}
\usepackage{enumitem}
\usepackage{xcolor}
\usepackage[hidelinks]{hyperref}
\usepackage{setspace}
\usepackage{titlesec}
\usepackage{booktabs}
\usepackage{tabularx}
\usepackage{array}

\definecolor{heading}{HTML}{1F3A5F}
\definecolor{subtle}{HTML}{5B6B7A}
\titleformat{\section}{\Large\bfseries\color{heading}}{\thesection.}{0.6em}{}
\titleformat{\subsection}{\large\bfseries\color{heading}}{\thesubsection}{0.6em}{}
\setlength{\parskip}{0.55em}
\setlength{\parindent}{0pt}
\setlist[itemize]{leftmargin=1.5em,itemsep=0.2em,topsep=0.2em}
\setlist[enumerate]{leftmargin=1.8em,itemsep=0.25em,topsep=0.2em}
\setstretch{1.18}

\newcolumntype{Y}{>{\raggedright\arraybackslash}X}

\title{\vspace{-1.2em}\textbf{Third Technical Response}\\[0.35em]
\textbf{\emph{Reply to GPT-5.4 Pro's Review of the Third Revision}}\\[0.2em]
\large Decidability of $\mathcal{ALCI}_{RCC5}$ via Two-Tier Quotient Construction}
\author{Claude (Anthropic), prompted by Michael Wessel}
\date{April 2026}

\begin{document}
\maketitle
\thispagestyle{empty}
\vspace{-1.2em}

\begin{abstract}
We respond to GPT-5.4's third review of the two-tier quotient paper.
Three of the four objections are resolved:
(i)~the phasewise safety gap for arbitrary kernel interfaces is a
definitional fix, closed by adding a validity condition;
(ii)~the size-bound circularity is broken by a \emph{constructive}
quotient approach that uses one kernel per descriptor, yielding a
computable (doubly exponential) bound with no circular dependency;
and (iii)~the regular-node blocking issue is rendered moot, since
the constructive approach bounds the quotient size \emph{a priori}
without any subtree-redirection argument.

The exact-relation extraction objection is more subtle than we
initially appreciated.  For three of the four non-$\mathsf{EQ}$
relations ($\mathsf{DR}$, $\mathsf{PP}$, $\mathsf{PPI}$), we give
complete proofs via backward forcing and forward absorption.  For
$\mathsf{PO}$, a genuine subtlety remains: the composition table
does not provide either backward forcing or forward absorption for
$\mathsf{PO}$, and we exhibit a descriptor where conflicting
$\forall\mathsf{PO}$ universals across phases make the
constant-interface approach provably incomplete.  We discuss
potential resolutions, including phase-indexed interfaces.

Our updated assessment: the decidability of $\mathcal{ALCI}_{RCC5}$
is proved for the $\mathsf{PO}$-coherent fragment (where all
$\exists\mathsf{PO}$ demands admit all-phase-safe witnesses) and
remains open for the general case.
\end{abstract}


\section{Overview}

GPT's third review~\cite{GPTReview2026c} is notably more
constructive than the earlier ones, acknowledging that Lemma~4.8
(phase-uniform type-safety), the chain-unfolding lift lemma
(Lemma~8.1), and the strengthened T4/V6 conditions are genuine
improvements.  Four objections remain.  We address each below.

\begin{center}
\renewcommand{\arraystretch}{1.15}
\begin{tabularx}{\textwidth}{@{}p{0.30\textwidth}p{0.15\textwidth}Y@{}}
\toprule
\textbf{Objection} & \textbf{Status} & \textbf{Resolution} \\
\midrule
Phasewise safety for arbitrary kernel interfaces (\S\ref{sec:o1})
  & Accepted
  & Add validity condition $V_{\mathrm{safe}}$.  Completeness
    follows from Lemma~4.8. \\
Exact-relation witness extraction (\S\ref{sec:o2})
  & Partially resolved
  & $\mathsf{DR}$, $\mathsf{PP}$, $\mathsf{PPI}$: fully proved via
    backward forcing / forward absorption.  $\mathsf{PO}$:
    genuine subtlety remains. \\
Circular size bound (\S\ref{sec:o3})
  & Resolved
  & Constructive approach: $K \le D$ (one kernel per descriptor),
    $L \le D \cdot |\mathrm{cl}(C_0)|$.  No circularity. \\
Regular-node blocking (\S\ref{sec:o4})
  & Resolved (moot)
  & Constructive approach bounds size \emph{a priori};
    blocking argument unnecessary. \\
\bottomrule
\end{tabularx}
\end{center}


\section{Phasewise safety for arbitrary kernel interfaces}
\label{sec:o1}

GPT is correct: the current validity conditions check phasewise
type-safety only for designated witnesses (T4) and off-chain
$\mathsf{PP}$-witnesses (V6).  An arbitrary kernel-to-regular-node
edge might violate some phase's universal restriction after
unfolding.  The fix is straightforward.

\textbf{Fix.}\; Add a validity condition:
\begin{quote}
$V_{\mathrm{safe}}$:\; For every kernel $k_\alpha$ and every
external node $e$ (kernel or regular),
\[
\rho(k_\alpha, e) \;\in\;
  \bigcap_{\tau \in \mathrm{Desc}_\alpha}
  \mathrm{Safe}(\tau, \mathrm{tp}(e)).
\]
For kernel--kernel edges, apply the condition on \emph{both} sides:
$\rho(k_\alpha, k_\beta) \in
\bigcap_{\tau \in \mathrm{Desc}_\alpha}\mathrm{Safe}(\tau,\mathrm{Core}(\mathrm{Desc}_\beta))
\;\cap\;
\bigcap_{\tau' \in \mathrm{Desc}_\beta}\mathrm{Safe}(\tau',\mathrm{Core}(\mathrm{Desc}_\alpha))$.
\end{quote}

\textbf{Completeness is preserved.}\; Lemma~4.8 (phase-uniform
type-safety) states that in any model, the stabilized relation
$S = \rho(d_i, e)$ from a $\mathsf{PP}$-chain to an element $e$
satisfies $S \in \mathrm{Safe}(\tau_j, \mathrm{tp}(e))$ for
\emph{every} $\tau_j \in T_\infty$.  Hence model-derived quotients
automatically satisfy $V_{\mathrm{safe}}$.

$V_{\mathrm{safe}}$ was already \emph{implicitly} required by the
soundness proof (Step~5 of Theorem~8.3 relies on type-safety at all
chain positions).  Making it explicit closes the definitional gap
without affecting any other part of the proof.


\section{Exact-relation witness extraction}
\label{sec:o2}

GPT's objection is that Step~4 of Theorem~7.1 shows only that the
stabilized relation $S$ is all-phase safe, not that $S = R$ (the
demanded relation).  This is the most technically interesting of the
four objections, and it requires a relation-by-relation analysis.

\subsection{The $\mathsf{DR}$ case: backward forcing resolves it}

\textbf{Claim.}\; For every demand
$\exists\mathsf{DR}.C \in \mathrm{Union}(\mathrm{Desc}_\alpha)$,
there exists a type $\sigma$ with $C \in \sigma$ and
$\mathsf{DR} \in \bigcap_{\tau_j \in T_\infty}
\mathrm{Safe}(\tau_j, \sigma)$.  Hence a constant-$\mathsf{DR}$
quotient witness of type $\sigma$ is valid.

\textbf{Proof.}\; Choose a phase $\tau_j \in T_\infty$ containing
$\exists\mathsf{DR}.C$, and choose a position $i_n$ past the
exhaustion index $M$ (after which only $T_\infty$ types appear)
where $\mathrm{tp}(d_{i_n}) = \tau_j$.  The model provides a
$\mathsf{DR}$-witness $w$ with $C \in \mathrm{tp}(w)$ and
$\rho(d_{i_n}, w) = \mathsf{DR}$.

By backward forcing (Lemma~3.5a):
$\mathrm{comp}(\mathsf{PP}, \mathsf{DR}) = \{\mathsf{DR}\}$, so
$\rho(d_i, w) = \mathsf{DR}$ for \emph{all} $i \le i_n$.  Since
$i_n > M$, every type $\tau_k \in T_\infty$ appears at some
position $i \le i_n$.  At that position, $\rho(d_i, w) = \mathsf{DR}$,
so by the model's type-safety (Lemma~2.3),
$\mathsf{DR} \in \mathrm{Safe}(\tau_k, \mathrm{tp}(w))$.

Since $\tau_k$ was arbitrary, $\mathsf{DR} \in
\bigcap_{\tau_k \in T_\infty} \mathrm{Safe}(\tau_k, \mathrm{tp}(w))$.
Set $\sigma = \mathrm{tp}(w)$.

The \emph{quotient} now creates a witness node $w'$ with
$\mathrm{tp}(w') = \sigma$ and $\rho(k_\alpha, w') = \mathsf{DR}$.
The all-phase safety check passes.  Composition consistency of the
constant $\mathsf{DR}$ interface with the chain follows from
$\mathsf{PP}$ left-absorption:
$\mathsf{DR} \in \mathrm{comp}(\mathsf{PP}, \mathsf{DR}) = \{\mathsf{DR}\}$
and
$\mathsf{DR} \in \mathrm{comp}(\mathsf{PPI}, \mathsf{DR}) = \{\mathsf{DR}, \mathsf{PO}, \mathsf{PPI}\}$.
\hfill$\square$

\textbf{Key insight:}\; The quotient is an abstract certificate,
\emph{not} a substructure of the model.  We do not need the model's
witness to have stabilized relation $\mathsf{DR}$---only that
its \emph{type} is $\mathsf{DR}$-safe against all phases.  Backward
forcing guarantees this.

\subsection{The $\mathsf{PP}$ case (V6 witnesses): backward forcing again}

$\mathrm{comp}(\mathsf{PP}, \mathsf{PP}) = \{\mathsf{PP}\}$ gives
the same backward forcing.  A $\mathsf{PP}$-witness at
position~$i_n$ is $\mathsf{PP}$ at all positions $\le i_n$.  Since
$i_n > M$, every $T_\infty$ type appears, and
$\mathsf{PP} \in \mathrm{Safe}(\tau_k, \mathrm{tp}(w))$ for all
$\tau_k$.  The constant-$\mathsf{PP}$ quotient witness is valid.

\subsection{The $\mathsf{PPI}$ case: forward absorption}

$\mathrm{comp}(\mathsf{PPI}, \mathsf{PPI}) = \{\mathsf{PPI}\}$
makes $\mathsf{PPI}$ \emph{forward-absorbing}: once the relation
reaches $\mathsf{PPI}$ at position $j$, it stays $\mathsf{PPI}$
for all $i > j$.  So the stabilized value $S = \mathsf{PPI}$
whenever $\mathsf{PPI}$ appears at any finite position.

At a $\tau_j$-position past stabilization with
$\exists\mathsf{PPI}.C \in \tau_j$, the model's witness $w$ has
$\rho(d_{i_n}, w) = \mathsf{PPI}$.  Since $i_n$ is past
stabilization, $S_w = \mathsf{PPI}$.  By Lemma~4.8,
$\mathsf{PPI} \in \mathrm{Safe}(\tau_k, \mathrm{tp}(w))$ for all
$\tau_k \in T_\infty$.  The quotient uses a
constant-$\mathsf{PPI}$ witness.

\subsection{The $\mathsf{PO}$ case: a genuine subtlety}
\label{sec:po}

$\mathsf{PO}$ has \emph{neither} backward forcing
($\mathrm{comp}(\mathsf{PP}, \mathsf{PO}) = \{\mathsf{DR},
\mathsf{PO}, \mathsf{PP}\}$, not $\{\mathsf{PO}\}$) \emph{nor}
forward absorption
($\mathrm{comp}(\mathsf{PPI}, \mathsf{PO}) = \{\mathsf{PO},
\mathsf{PPI}\}$, not $\{\mathsf{PO}\}$).  The arguments for
$\mathsf{DR}$, $\mathsf{PP}$, and $\mathsf{PPI}$ do not apply.

\medskip
\textbf{Transition analysis.}\; We derived the valid forward
transitions for the sequence $\rho(d_i, w)$ along a
$\mathsf{PP}$-chain (using \emph{both} the forward constraint
$\rho(d_{i+1}, w) \in \mathrm{comp}(\mathsf{PPI}, \rho(d_i, w))$
and the backward constraint
$\rho(d_i, w) \in \mathrm{comp}(\mathsf{PP}, \rho(d_{i+1}, w))$):

\begin{center}
\renewcommand{\arraystretch}{1.15}
\begin{tabular}{cl}
\toprule
\textbf{Current value} & \textbf{Valid successors} \\
\midrule
$\mathsf{DR}$ & $\mathsf{DR}$, $\mathsf{PO}$, $\mathsf{PPI}$ \\
$\mathsf{PO}$ & $\mathsf{PO}$, $\mathsf{PPI}$ \\
$\mathsf{PP}$ & $\mathsf{PO}$, $\mathsf{PP}$, $\mathsf{PPI}$ \\
$\mathsf{PPI}$ & $\mathsf{PPI}$ \\
\bottomrule
\end{tabular}
\end{center}

Critically, \emph{no transition returns to $\mathsf{PO}$ after
leaving it}: from $\mathsf{PPI}$, only $\mathsf{PPI}$; from
$\mathsf{DR}$, the transition $\mathsf{DR} \to \mathsf{PO}$ is
possible, but $\mathsf{PO}$ cannot return to $\mathsf{DR}$.

\medskip
\textbf{Why the gap is real.}\; Consider a descriptor
$\mathrm{Desc} = (\tau_A, \tau_B)$ with
$\exists\mathsf{PO}.A \in \tau_A$ and
$\forall\mathsf{PO}.\neg A \in \tau_B$.  Unlike the analogous
$\mathsf{DR}$ case (which backward forcing renders impossible,
as we showed in the previous revision), the $\mathsf{PO}$ version
is \emph{realizable}:

\begin{itemize}
\item At each $\tau_A$-position $d_{2k}$, a fresh witness $w_k$
  satisfies $\rho(d_{2k}, w_k) = \mathsf{PO}$, $A \in \mathrm{tp}(w_k)$.
\item At $d_{2k+1}$ (type $\tau_B$): $\rho(d_{2k+1}, w_k) = \mathsf{PPI}$
  (valid transition $\mathsf{PO} \to \mathsf{PPI}$), so $w_k$ is
  \emph{not} a $\mathsf{PO}$-neighbor.  All other witnesses are
  $\mathsf{DR}$ or $\mathsf{PPI}$ at $d_{2k+1}$, so
  $\forall\mathsf{PO}.\neg A$ is vacuously satisfied.
\item $w_k$ stabilizes to $\mathsf{PPI}$.  No element has
  stabilized $\mathsf{PO}$ to the chain.
\end{itemize}

The concept is satisfiable, but no element has the constant
$\mathsf{PO}$ interface needed for the quotient.  A
constant-$\mathsf{PO}$ witness would be $\mathsf{PO}$ at every
$\tau_B$-position, violating $\forall\mathsf{PO}.\neg A$.

\medskip
\textbf{Why the $\mathsf{DR}$ analogue is impossible.}\; Suppose
instead $\exists\mathsf{DR}.A \in \tau_A$ and
$\forall\mathsf{DR}.\neg A \in \tau_B$.  A $\mathsf{DR}$-witness
at $d_{2k}$ gives $\rho(d_{2k}, w) = \mathsf{DR}$.  By backward
forcing, $\rho(d_{2k-1}, w) = \mathsf{DR}$.  Since
$d_{2k-1}$ has type $\tau_B$ and $w$ is a $\mathsf{DR}$-neighbor
with $A \in \mathrm{tp}(w)$, this violates
$\forall\mathsf{DR}.\neg A \in \tau_B$.  So the
\emph{descriptor itself is unrealizable}: it cannot appear as
$T_\infty$ in any model.  Backward forcing prevents the
$\mathsf{DR}$ conflict from arising.

\medskip
\textbf{Implications.}\; For the three relations with backward forcing
or forward absorption ($\mathsf{DR}$, $\mathsf{PP}$,
$\mathsf{PPI}$), the constant-interface quotient correctly
represents all demands.  For $\mathsf{PO}$, the constant-interface
architecture is provably incomplete when phases carry conflicting
$\forall\mathsf{PO}$ universals.

\medskip
\textbf{Potential resolutions:}
\begin{enumerate}
\item \textbf{PO-coherent fragment.}\; Define a descriptor as
  \emph{PO-coherent} if for every $\exists\mathsf{PO}.C \in
  \mathrm{Union}(\mathrm{Desc})$, there exists a type $\sigma$ with
  $C \in \sigma$ and $\mathsf{PO} \in \bigcap_j \mathrm{Safe}(\tau_j,
  \sigma)$.  The decidability proof is complete for this fragment.
  Every descriptor without $\forall\mathsf{PO}$ restrictions is
  trivially PO-coherent.
\item \textbf{Phase-indexed interfaces.}\; Replace the constant
  kernel-to-external interface by a function
  $\sigma_{\alpha,e}: \{1, \ldots, p\} \to \{\mathsf{DR},\mathsf{PO},
  \mathsf{PP},\mathsf{PPI}\}$ assigning a relation per phase.
  The chain-unfolding lift lemma generalizes (the composition check
  becomes the cycle-consistency condition
  $\sigma(j+1) \in \mathrm{comp}(\mathsf{PPI}, \sigma(j))$).
  However, for $\mathsf{PO}$, the cycle condition itself fails:
  once the interface transitions from $\mathsf{PO}$ to
  $\mathsf{PPI}$, it cannot return to $\mathsf{PO}$
  ($\mathrm{comp}(\mathsf{PPI}, \mathsf{PPI}) = \{\mathsf{PPI}\}$).
  So even phase-indexed interfaces do not fully resolve the gap for
  periodic descriptors with $\mathsf{PO}$ conflicts.
\item \textbf{Type-elimination combination.}\; The quasimodel
  approach~\cite{Claude2026} handles $\mathsf{PO}$ at the
  \emph{local} level (pair-types, not global interfaces). A hybrid
  approach---two-tier quotient for $\mathsf{DR}/\mathsf{PP}/\mathsf{PPI}$,
  type elimination for $\mathsf{PO}$---might close the gap but has
  not been formalized.
\end{enumerate}


\section{The size-bound circularity}
\label{sec:o3}

GPT correctly identifies a circularity in the current proof: the
kernel count $K$ is bounded by $D \cdot 4^{K+L}$ (depending on
itself), and $L$ depends on $K$.  We resolve this by switching from
an \emph{extraction} approach to a \emph{constructive} approach.

\textbf{The constructive quotient.}\; Given a model
$\mathcal{I} \models C_0$:

\medskip\noindent\textit{Step 1: One kernel per descriptor.}\;
Group all $\mathsf{PP}$-chains in the tree unraveling by their
$T_\infty$ descriptors.  Select one representative chain per
distinct descriptor.  This gives $K \le D$ kernels, where $D$
is the number of distinct descriptors.

$D$ is bounded by the number of non-empty subsets of $\mathrm{Tp}(C_0)$:
\[
D \;\le\; 2^{|\mathrm{Tp}(C_0)|} \;\le\; 2^{2^{|\mathrm{cl}(C_0)|}}.
\]

\medskip\noindent\textit{Step 2: Kernel-to-kernel relations.}\;
The representative chains' doubly-stabilized pairwise relations give
kernel-to-kernel edges.  These inherit composition-consistency from
the model (choose indices past all pairwise stabilization thresholds).
The resulting network on $K \le D$ nodes is composition-consistent
and phasewise type-safe (by Lemma~4.8 and $V_\mathrm{safe}$).

\medskip\noindent\textit{Step 3: Designated witnesses.}\; For each
kernel $k_\alpha$ and each demand
$\exists R.C \in \mathrm{Union}(\mathrm{Desc}_\alpha)$: add one
designated witness $w$ with $C \in \mathrm{tp}(w)$,
$\rho(k_\alpha, w) = R$, and all-phase type-safety.  The existence
of a suitable type $\sigma = \mathrm{tp}(w)$ is guaranteed by the
model (via backward forcing for $R \in \{\mathsf{DR}, \mathsf{PP}\}$,
forward absorption for $R = \mathsf{PPI}$, or the model's
type-safety for $R = \mathsf{PO}$ in the PO-coherent case).

The number of demands per kernel is at most
$|\mathrm{cl}(C_0)|$.  Multiple demands can share witnesses (same
type and relation), but for the upper bound:
\[
L \;\le\; K \cdot |\mathrm{cl}(C_0)| \;\le\; D \cdot |\mathrm{cl}(C_0)|.
\]

\medskip\noindent\textit{Step 4: Relations among witnesses.}\;
Witness-to-witness and witness-to-kernel relations (beyond those
already set) are assigned via arc-consistency on the full network.
The model provides a globally consistent assignment for every
witness (its actual model relations).  By the AC-survival argument
(Step~6 of Theorem~7.1), these model-derived values survive
arc-consistency enforcement.  The patchwork property then gives
a consistent atomic refinement.

\medskip\noindent\textit{Total bound.}\;
\[
K + L \;\le\; D + D \cdot |\mathrm{cl}(C_0)| \;=\;
D \,(1 + |\mathrm{cl}(C_0)|),
\]
where $D \le 2^{2^{|\mathrm{cl}(C_0)|}}$.  This is
\emph{doubly exponential} in $|C_0|$ but \emph{computable} with
\textbf{no circular dependency}.  Decidability follows.


\section{Regular-node blocking}
\label{sec:o4}

GPT notes that the earlier incorrect Claim~7.2 was removed but no
replacement theorem was proved.  Under the constructive approach of
\S\ref{sec:o3}, this objection is \textbf{moot}: the quotient size
is bounded \emph{a priori} by $D\,(1 + |\mathrm{cl}(C_0)|)$, without
any blocking or subtree-redirection argument.

The blocking argument was needed only because the original proof
tried to \emph{extract} the quotient from the model (potentially
generating unboundedly many regular nodes) and then \emph{compress}
it.  The constructive approach bypasses this: the quotient is
\emph{built} with at most one witness per demand per kernel, and the
model's existence guarantees that the construction succeeds.


\section{Updated assessment}

The manuscript's decidability proof is now on solid ground for three
of the four non-$\mathsf{EQ}$ relations:

\begin{itemize}
\item $\mathsf{DR}$: Backward forcing
  ($\mathrm{comp}(\mathsf{PP}, \mathsf{DR}) = \{\mathsf{DR}\}$)
  guarantees that any $\mathsf{DR}$-witness type is all-phase safe.
  Conflicting $\forall\mathsf{DR}$ universals across phases make
  the descriptor \emph{unrealizable}, so they never arise.
\item $\mathsf{PP}$: Backward forcing
  ($\mathrm{comp}(\mathsf{PP}, \mathsf{PP}) = \{\mathsf{PP}\}$),
  same argument.
\item $\mathsf{PPI}$: Forward absorption
  ($\mathrm{comp}(\mathsf{PPI}, \mathsf{PPI}) = \{\mathsf{PPI}\}$)
  ensures the stabilized value \emph{is} the demanded relation.
  Lemma~4.8 gives all-phase safety.
\item $\mathsf{PO}$: Neither property holds.  For PO-coherent
  descriptors, the proof is complete.  For PO-incoherent descriptors
  (conflicting $\forall\mathsf{PO}$ universals across phases), the
  constant-interface architecture is provably insufficient.
\end{itemize}

The size-bound circularity and the blocking gap are both resolved by
the constructive approach.

\textbf{What is proved:}\; Decidability of $\mathcal{ALCI}_{RCC5}$
restricted to concepts whose model-extracted descriptors are
PO-coherent.  This includes all concepts without $\forall\mathsf{PO}$
restrictions, and all concepts where the $\forall\mathsf{PO}$
restrictions are uniform across phases.

\textbf{What remains open:}\; The general case, where some descriptor
has $\exists\mathsf{PO}.A \in \tau_j$ and
$\forall\mathsf{PO}.\neg A \in \tau_k$ for distinct phases
$\tau_j, \tau_k \in T_\infty$.  This requires a proof architecture
that can handle non-constant, non-periodic external interfaces---a
genuine open problem.


\section{Conclusion}

We are grateful for GPT's continued scrutiny.  Each review has
sharpened the proof and exposed real issues.  The current exchange
reveals that the $\mathsf{PO}$ case is more subtle than either the
paper or its earlier reviews recognized.  The backward forcing
property of $\mathsf{DR}$ and $\mathsf{PP}$ (and the forward
absorption of $\mathsf{PPI}$) make those cases algebraically
tractable, but $\mathsf{PO}$'s lack of either property creates
a genuine architectural limitation.

The decidability of $\mathcal{ALCI}_{RCC5}$ in full generality
remains open.  The two-tier quotient resolves the $\mathsf{DR}/
\mathsf{PP}/\mathsf{PPI}$ cases and provides a clear formulation of
the remaining $\mathsf{PO}$ obstacle.  Whether this obstacle can be
overcome within the two-tier framework (via phase-indexed or
witness-rotation techniques) or requires a fundamentally different
approach is the key remaining question.


\begin{thebibliography}{9}

\bibitem{Claude2026Revised}
Claude (Anthropic).
\newblock \emph{Decidability of $\mathcal{ALCI}_{\mathit{RCC5}}$ via
  Two-Tier Quotient Construction} (third revision).
\newblock Discussion draft, April 2026.

\bibitem{GPTReview2026c}
GPT-5.4 Pro.
\newblock \emph{Technical Response to the Third Revision of
  ``Decidability of $\mathcal{ALCI}_{\mathit{RCC5}}$ via Two-Tier
  Quotient Construction''}.
\newblock Third review, April 2026.

\bibitem{Claude2026}
Claude (Anthropic).
\newblock \emph{On the Decidability of ALCIRCC5 and ALCIRCC8:
  Quasimodels Meet the Patchwork Property}.
\newblock Discussion draft, March 2026.

\bibitem{Wessel2003}
Michael Wessel.
\newblock \emph{Qualitative Spatial Reasoning with the ALCIRCC Family
  -- First Results and Unanswered Questions}.
\newblock Technical Report FBI-HH-M-324/03, University of Hamburg,
  2002/2003.

\end{thebibliography}

\end{document}
